% !TeX program = lualatex
% !TeX encoding = utf8
% !TeX spellcheck = uk_UA
% !TeX root =../PyscfBook.tex

%=========================================================
\Opensolutionfile{answer}[\currfilebase/\currfilebase-Answers]
\chapter{Коливальні властивості молекул}\label{\currfilebase}
%=========================================================

Коливальна спектроскопія --- один із найпотужніших методів ідентифікації
молекул та вивчення їх структури. Інфрачервоні (ІЧ) та Раман спектри
надають інформацію про коливальні моди, силові константи зв'язків,
симетрію молекули, а також --- через гармонійні частоти ---
\emph{нульову енергію коливань} (ZPE), що необхідна для точного
розрахунку ентальпії, ентропії та вільної енергії в термохімії.

%\textbf{Зв’язок між коливаннями та спектроскопією.}
%Коливальні властивості молекул --- це, насамперед, набір нормальних мод,
%кожна з яких характеризується власною частотою, амплітудою та типом руху
%(розтягування, згинання, скручування тощо). Ці моди не можна побачити
%«неозброєним оком», але їх можна \emph{збудити} і \emph{зареєструвати}
%за допомогою електромагнітного випромінювання. Саме тут вступає в дію
%\textbf{коливальна спектроскопія}: вона перетворює абстрактні механічні
%коливання на вимірювані спектральні сигнали --- піки в ІЧ- або Раман-спектрах.

\textbf{Частоти в різних одиницях:}
\begin{center}
    \begin{tblr}{
        colspec = {Q[l] Q[r] X[c]},
        row{1} = {c, m, font=\bfseries},
        hline{1,2,Z} = {solid},
            }
        Одиниця   & Множник                & Використання             \\
        см$^{-1}$ & 1                      & ІЧ спектроскопія         \\
        Гц        & $2.998 \times 10^{10}$ & СІ одиниці               \\
        еВ        & $1.240 \times 10^{-4}$ & Електронна спектроскопія \\
        Хартрі    & $4.556 \times 10^{-6}$ & Атомні одиниці           \\
    \end{tblr}
\end{center}



%% --------------------------------------------------------
\section{Гессіан та нормальні моди}
%% --------------------------------------------------------

Коливання молекули відбуваються поблизу рівноважної геометрії $\mathbf{R}_0$.
Для малих відхилень потенціальну енергію апроксимують квадратичною формою:
\[
    E(\mathbf{R}) \approx E(\mathbf{R}_0) +
    \frac{1}{2} \Delta\mathbf{R}^T \, \mathbf{H} \, \Delta\mathbf{R},
\]
де $\mathbf{H}$ --- \textbf{гесіан} (матриця других похідних):
\[
    H_{ij} = \left. \frac{\partial^2 E}{\partial R_i \partial R_j} \right|_{\mathbf{R}_0}.
\]
Для $N$ атомів гесіан має розмір $3N \times 3N$ і повністю описує гармонічні коливання.

%% --------------------------------------------------------
\subsection{Матриця силових констант та діагоналізація}
%% --------------------------------------------------------

Елементи $H_{ij}$ --- це \textbf{силові константи}:
діагональні --- жорсткість окремого руху,
недіагональні --- зв’язок між різними рухами
(наприклад, розтяг \ce{C-H} та згин \ce{H-C-O}).

У декартових чи внутрішніх координатах коливання \emph{пов’язані}.
\textbf{Діагоналізація гесіану} дає:
\begin{itemize}
    \item \textit{власні значення} $\omega_k^2$ --- квадрати частот;
    \item \textit{власні вектори} $\mathbf{L}_k$ --- \emph{нормальні моди} (незалежні коливання).
\end{itemize}

Це роблять \textbf{після оптимізації геометрії}, щоб:
\begin{itemize}
    \item підтвердити мінімум ($\omega_k^2 > 0$);
    \item обчислити \textbf{ZPE}: $E_{\text{ZPE}} = \frac{1}{2} h \sum \omega_k$;
    \item знайти внесок у термохімію ($H$, $S$, $G$);
    \item змоделювати ІЧ/Раман-спектри (піки при $\omega_k$).
\end{itemize}

Отже, з гесіану через нормальні моди отримують повну картину коливальних властивостей молекули.

%% --------------------------------------------------------
\subsection{Нормальні коливання}
%% --------------------------------------------------------

Для переходу від гесіану $\mathbf{H}$ у декартових координатах до фізично осмислених частот враховують маси атомів.
Вводять \textbf{мас-зважений гесіан}:
\[
    \mathbf{F} = \mathbf{M}^{-1/2} \mathbf{H} \mathbf{M}^{-1/2},
\]
де $\mathbf{M}$ --- діагональна матриця атомних мас.

Далі розв'язують задачу на власні значення:
\[
    \mathbf{F} \mathbf{Q} = \omega^2 \mathbf{Q},
\]
де $\omega$ --- коливальні частоти, а $\mathbf{Q}$ --- вектори нормальних координат.

Отримані власні значення $\omega_i^2$ відповідають:
\begin{itemize}
    \item $3N - 6$ коливальним модам для нелінійної молекули;
    \item $3N - 5$ --- для лінійної;
    \item 3 трансляційним і 3 (або 2) обертальним модам з $\omega = 0$.
\end{itemize}

\paragraph{Зв'язок координат $Q_k$ та $\mathbf{R}$}

Нормальні координати $Q_k$ є лінійними комбінаціями атомних зміщень $\Delta \mathbf{R}$:
\[
    Q_k = \sum_{i=1}^{3N} \sqrt{m_i}\, l_{ik}\, \Delta R_i
\]
або у матричній формі:
\[
    \mathbf{Q} = \mathbf{L}^T \mathbf{M}^{1/2} \Delta \mathbf{R},
\]
де:
\begin{itemize}
    \item $\mathbf{M}$ — діагональна матриця мас атомів,
    \item $\mathbf{L}$ — матриця власних векторів (нормальних мод),
    \item $\Delta \mathbf{R} = \mathbf{R} - \mathbf{R}_0$ — відхилення координат від рівноваги.
\end{itemize}

Звідси обернене перетворення:
\[
    \Delta \mathbf{R} = \mathbf{M}^{-1/2} \mathbf{L} \, Q_k.
\]

Саме цю формулу використовують у чисельному диференціюванні:
\[
    \mathbf{R}^{(\pm)} = \mathbf{R}_0 \pm h\, \mathbf{M}^{-1/2} \mathbf{L}_k,
\]
що зв’язує крок $h$ у просторі нормальних координат $Q_k$ із фізичними зміщеннями атомів $\mathbf{R}$.


%% --------------------------------------------------------
\subsubsection{Масштабування частот}
%% --------------------------------------------------------

Розраховані гармонічні частоти $\nu_{\text{гарм}}$ систематично завищені відносно експериментальних $\nu_{\text{експ}}$ через спрощення моделі потенціальної поверхні та обмеження методів:

\begin{itemize}
    \item \textbf{Гармонічне наближення:} реальний потенціал ангармонічний, тому енергія коливань і частоти менші.
    \item \textbf{Нестача електронної кореляції:} метод Хартрі–Фока завищує жорсткість зв’язків.
    \item \textbf{Обмежений базис:} малий базис перебільшує градієнти енергії.
\end{itemize}

{Масштабування $\lambda$ емпірично компенсує ангармонічність і методичні похибки. Емпіричне виправлення:
\[
    \nu_{\text{корект}} = \lambda \cdot \nu_{\text{гарм}}
\]

\textbf{Типові масштабувальні фактори:}

\begin{center}
    \begin{tblr}{
        colspec = {X[l,2] X[r] X[l,2]},
        row{1} = {font=\bfseries},
        hline{1,2,Z} = {solid},
            }
        Метод/базис         & $\lambda$ & Коментар               \\
        HF/6-31G*           & 0.8929    & Сильно завищує         \\
        B3LYP/6-31G*        & 0.9613    & Універсальний          \\
        B3LYP/6-311+G(2d,p) & 0.9679    & Кращий базис           \\
        MP2/6-31G*          & 0.9434    & Для точних розрахунків \\
    \end{tblr}
\end{center}

\textit{Джерело: NIST Computational Chemistry Comparison and Benchmark Database}



%% --------------------------------------------------------
\subsection{Обчислення гесіану}
%% --------------------------------------------------------

PySCF може обчислювати гесіан аналітично або чисельно. Програма в файлі \inputcode{h2o_hessian.py} обчислює матрицю Гесіану для молекули води (\ce{H2O}) двома методами: чисельним та аналітичним.

\begin{commentbox}
У програмному пакеті \texttt{PySCF} гесіан зберігається не як звичайна
двовимірна матриця розмірності $3N \times 3N$, а у вигляді чотиривимірного
тензора
\[
H_{AB,ij} = \frac{\partial^2 E}{\partial R_{A,i} \partial R_{B,j}},
\]
де $A,B$~— індекси атомів ($1,\dots,N$), а $i,j \in \{x,y,z\}$~—
напрямки декартових координат.
Відповідно, внутрішня структура масиву має вигляд
\[
H \in \mathbb{R}^{N \times N \times 3 \times 3}.
\]

Щоб отримати стандартну матричну форму $3N \times 3N$,
необхідно спершу переставити осі, розмістивши напрямки координат поруч із
відповідними атомами:
\[
H'_{Ai,Bj} = H_{AB,ij}.
\]
Це виконується операцією
\[
\texttt{H'} = \inlinecode{H.transpose(0,2,1,3).reshape(3N,3N)}.
\]
Така трансформація формує коректну блочну структуру з блоками $3\times3$,
яка відповідає фізичному гесіану молекули у декартових координатах.
\end{commentbox}

\subsubsection{Чисельний Гесіан}

Чисельний Гесіан обчислюється через скінченні різниці методом центральних різниць:

\begin{enumerate}
    \item \textbf{Зміщення координат:} Кожна координата $x_i$ кожного атома зміщується на величину $\pm h$ (крок диференціювання):
    \begin{equation}
        x_i^{(\pm)} = x_i \pm h
    \end{equation}
    де $h = 0.001$ Bohr (атомні одиниці).

    \item \textbf{Обчислення градієнтів:} Для кожного зміщення обчислюється градієнт енергії:
    \begin{equation}
        \nabla E(x_i^{(+)}) = \left( \frac{\partial E}{\partial x_1}, \frac{\partial E}{\partial x_2}, \ldots, \frac{\partial E}{\partial x_N} \right)_{x_i = x_i + h}
    \end{equation}
    \begin{equation}
        \nabla E(x_i^{(-)}) = \left( \frac{\partial E}{\partial x_1}, \frac{\partial E}{\partial x_2}, \ldots, \frac{\partial E}{\partial x_N} \right)_{x_i = x_i - h}
    \end{equation}

    \item \textbf{Обчислення Гесіану:} Елементи матриці Гесіану обчислюються як:
    \begin{equation}
        H_{ij} = \frac{\partial^2 E}{\partial x_i \partial x_j} = \frac{\nabla E_j(x_i^{(+)}) - \nabla E_j(x_i^{(-)})}{2h}
    \end{equation}
\end{enumerate}

\textbf{Формула центральних різниць:}
\begin{equation}
    \boxed{H_{ij} = \frac{g_j(x_i + h) - g_j(x_i - h)}{2h}}
\end{equation}
де $g_j = \frac{\partial E}{\partial x_j}$ --- компонента градієнта.

\subsubsection{Аналітичний Гесіан}

Аналітичний Гесіан обчислюється через точні другі похідні енергії, використовуючи модуль PySCF:

\begin{equation}
    H_{ij} = \frac{\partial^2 E}{\partial x_i \partial x_j}
\end{equation}

Метод використовує \inlinecode{pyscf.hessian.rhf.Hessian(mf).kernel()} для прямого обчислення всіх елементів матриці Гесіану без числового диференціювання.


\paragraph{Обчислювальна складність}

\begin{itemize}
    \item \textbf{Чисельний метод:} Потребує $2 \times 3N$ SCF розрахунків, де $N$ --- кількість атомів.
    \begin{itemize}
        \item Для H$_2$O ($N=3$): $2 \times 3 \times 3 = 18$ SCF розрахунків.
        \item Час виконання: $\sim 10 - 30$ секунд.
    \end{itemize}

    \item \textbf{Аналітичний метод:} Один розрахунок з аналітичними похідними.
    \begin{itemize}
        \item Час виконання: $\sim 1 - 5$ секунд.
        \item \textbf{Прискорення:} у $5 - 20$ разів швидше.
    \end{itemize}
\end{itemize}

\paragraph{Точність}

\begin{itemize}
    \item \textbf{Аналітичний метод:} Точніший, оскільки використовує точні похідні без похибок дискретизації

    \item \textbf{Чисельний метод:} Точність залежить від кроку $h$:
    \begin{equation}
        \text{Похибка} \sim O(h^2) + \frac{\epsilon}{h}
    \end{equation}
    де $\epsilon$ --- машинна точність

    \item \textbf{Типова різниця:} Максимальна абсолютна різниця $\sim 10^{-5}$--$10^{-7}$ Ha/Bohr$^2$.
\end{itemize}

\subsubsection{Коливальні частоти}

З матриці Гесіану можна отримати коливальні частоти молекули:

\begin{equation}
    \omega_i = \sqrt{\lambda_i} \times 5140.48,
\end{equation}
де $5140.48$~--- це числовий множник, який переводить частоти з атомних одиниць у~см$^{-1}$, а $\lambda_i$ --- власні значення масово-зваженого Гесіану:

\begin{equation}
    \tilde{H}_{ij} = \frac{H_{ij}}{\sqrt{m_i m_j}}
\end{equation}

\noindent\textbf{Очікувані частоти для H$_2$O:}
\begin{itemize}
    \item Перші 6 мод: $\approx 0$ см$^{-1}$ (трансляції та обертання)
    \item Симетричне валентне коливання: $\approx 3600$ см$^{-1}$
    \item Деформаційне коливання: $\approx 1600$ см$^{-1}$
    \item Асиметричне валентне коливання: $\approx 3800$ см$^{-1}$
\end{itemize}

Для практичних розрахунків \textbf{аналітичний метод є кращим вибором}, оскільки він значно швидший та точніший. Чисельний метод корисний для перевірки результатів або коли аналітичні похідні недоступні.

%% --------------------------------------------------------
\section{Інфрачервоні (ІЧ) спектри}
%% --------------------------------------------------------

\textbf{ІЧ-спектроскопія} досліджує взаємодію молекул з інфрачервоним випромінюванням.
Поглинання відбувається тоді, коли частота випромінювання збігається з частотою
внутрішнього коливання молекули, тобто виконується умова резонансу:
\[
    h\nu = \Delta E_{\text{колив.}}
\]

Ці переходи відповідають змінам \emph{коливальних} рівнів енергії.
Типові частоти для молекулярних коливань лежать у межах
\[
    400 - 4000~\text{см}^{-1},
\]
що відповідає довжинам хвиль у мікрометровому діапазоні ($2.5 - 25~\mu$m).

ІЧ-спектр відображає лише ті \textit{коливальні моди молекули}, які супроводжуються \emph{зміною її дипольного моменту}.

%% --------------------------------------------------------
\subsection{Інтенсивності ІЧ-поглинання}
%% --------------------------------------------------------

У спектроскопії інтенсивність смуги --- це \textbf{інтегральна поглинальна здатність}, тобто
площа під піком у спектрі.
Її фізичний зміст --- це міра того, \emph{наскільки сильно молекула поглинає ІЧ-випромінювання}
при переході $v=0 \to 1$.

Із теорії взаємодії електромагнітного поля з речовиною відомо, що інтенсивність $A_i$ для $i$-го коливального переходу визначається як:
\[
    A_i = \frac{N_A \pi}{3 c^2}
        \left| \frac{\partial \boldsymbol{\mu}}{\partial Q_i} \right|^2,
\]
де $\tfrac{\partial \boldsymbol{\mu}}{\partial Q_i}$ --- похідна дипольного моменту за нормальною координатою
$Q_i$ у рівноважному положенні, $N_A$ --- число Авогадро, а $c$ --- швидкість світла.

\paragraph{Фізичний зміст.}
Величина $A_i$ має розмірність площі на моль:
\[
    [A_i] = \text{см}^2 \text{/моль},
\]
тобто її можна трактувати як \textbf{ефективну площу поглинання на один моль молекул} речовини.
Чим більша ця площа --- тим інтенсивніше смуга в спектрі.

\paragraph{Зв’язок з експериментом.}
В експериментальній ІЧ-спектроскопії вимірюють \emph{оптичну густину} $A$ (безрозмірну величину),
яка пов’язана із концентрацією $c$ та довжиною кювети $l$ законом Бугера–Ламберта–Бера:
\[
    A = \varepsilon\, c\, l,
\]
де $\varepsilon$ --- молярний коефіцієнт поглинання.
Інтеграл від $\varepsilon(\tilde{\nu})$ за частотою $\tilde{\nu}$ визначає \textbf{молярну інтегральну інтенсивність смуги}:
\[
    I = \int \varepsilon(\tilde{\nu})\, d\tilde{\nu},
\]
яка має ті ж самі одиниці, що й $A_i$, тобто см$^2$/моль.

\paragraph{Одиниці км/моль.}
Щоб уникнути дуже малих чисел, цю площу традиційно виражають у \textbf{кілометрах на моль}:
\[
    1~\text{км/моль} = 10^5~\text{см}^2/\text{моль}.
\]
Тоді типові значення інтенсивностей для фундаментальних коливань молекул
лежать у межах 10–1000~км/моль, що зручно для порівняння та таблиць.

\paragraph{Числова форма.}
Після підстановки числових значень констант:
\[
    N_A = 6.022 \times 10^{23}~\text{моль}^{-1}, \quad
    c = 2.998 \times 10^{10}~\text{см/с},
\]
отримуємо:
\[
    \boxed{I = 974.8638 \times
    \left| \frac{\partial \boldsymbol{\mu}}{\partial Q} \right|^2_{\text{a.u.}} \text{ км/моль}.}
\]

\paragraph{Зв’язок між $A_i$, $I$ і коефіцієнтом $K$.}
Формально експериментальна інтегральна інтенсивність $I_i$
і теоретичне вираження $A_i$ описують одну й ту саму фізичну величину,
записану в різних одиницях:
\[
    I_i = K\,A_i.
\]
Коефіцієнт $K$ враховує перехід від атомних одиниць, у яких обчислюється
похідна $\tfrac{\partial \boldsymbol{\mu}}{\partial Q_i}$, до лабораторних одиниць (км/моль).
Підставивши фізичні сталі,
\[
    K = \frac{N_A \pi}{3 c^2} (e a_0)^2 \times 10^5,
\]
отримуємо числове значення
\[
    \boxed{K = 974.8638,}
\]
тому остаточний вираз для інтегральної інтенсивності має вигляд:
\[
    \boxed{I_i = 974.8638
    \left| \frac{\partial \boldsymbol{\mu}}{\partial Q_i} \right|^2_{\text{a.u.}}~
    [\text{км/моль}].}
\]



%% --------------------------------------------------------
\subsection{Класифікація за інтенсивністю}
%% --------------------------------------------------------

\begin{center}
\begin{tblr}{ll}
\toprule
\textbf{Інтенсивність (км/моль)} & \textbf{Класифікація} \\
\midrule
$I > 100$ & Сильна смуга \\
$10 < I \leq 100$ & Середня смуга \\
$I \leq 10$ & Слабка смуга \\
\bottomrule
\end{tblr}
\end{center}


%=========================================================
\section{Розрахунок ІЧ-інтенсивностей}
%=========================================================

ІЧ-інтенсивність нормальної моди пропорційна квадрату похідної дипольного моменту
за координатою цієї моди:
\[
I_i \sim \left| \frac{d\boldsymbol{\mu}}{dQ_i} \right|^2.
\]
На практиці цю похідну можна отримати двома різними шляхами, однак у

%---------------------------------------------------------
\subsection{Чисельне диференціювання}
%---------------------------------------------------------

У цьому підході координати атомів зсуваються вздовж $i$-ї нормальної моди
на малу величину~$\delta$, після чого обчислюються дипольні моменти
у двох геометріях --- \emph{«плюс»} і \emph{«мінус»}.
Похідна апроксимується центральною різницею:
\[
\left( \frac{d\boldsymbol{\mu}}{dQ_i} \right)_0 \approx
\frac{\boldsymbol{\mu}(Q_i+\delta) - \boldsymbol{\mu}(Q_i-\delta)}{2\delta}.
\]

В файлі \inputcode{h2o_ir_spectrum.py} наведено програму, яка  використовує \textbf{метод центральних різниць} для чисельного обчислення $\frac{d\boldsymbol{\mu}}{dQ_i}$:



\subsubsection{Алгоритм обчислення}

\begin{enumerate}
    \item \textbf{Отримання нормальної моди:}

    З діагоналізації масово-зваженого Гесіану отримуємо власний вектор $\mathbf{L}_k$ (нормальну моду):
    \begin{equation}
        \mathbf{L}_k = (l_{1x}, l_{1y}, l_{1z}, l_{2x}, l_{2y}, l_{2z}, \ldots, l_{Nx}, l_{Ny}, l_{Nz})^T
    \end{equation}
    де $N$ --- кількість атомів.

    \begin{minted}{python}
    freq_info = thermo.harmonic_analysis(mol, hess)
    normal_modes = freq_info['norm_mode']       # власні вектори

    L_k = normal_modes[:, k].reshape(natm, 3)
    \end{minted}

    Код \inlinecode{L\_k = normal\_modes[:, k].reshape(natm, 3)} відповідає $\mathbf{L}_k$ у формулі.

    \item \textbf{Зміщення координат:}

    Для чисельного диференціювання вибираємо крок $h$ та зміщуємо атоми вздовж нормальної моди:
    \begin{align}
        \mathbf{R}^{(+)} &= \mathbf{R}_0 + h \cdot \mathbf{L}_k \\
        \mathbf{R}^{(-)} &= \mathbf{R}_0 - h \cdot \mathbf{L}_k
    \end{align}

    \begin{minted}{python}
    h = 0.001  # Bohr
    R_0 = mol.atom_coords()
    R_pos = R_0 + h * L_k
    R_neg = R_0 - h * L_k
    \end{minted}

    \item \textbf{SCF-розрахунки:}

    Для кожної зміщеної геометрії виконуємо SCF-розрахунок, отримуємо енергію $E$ та дипольний момент $\boldsymbol{\mu}$:
    \begin{align}
        E^{(+)}, \boldsymbol{\mu}^{(+)} &= \text{SCF}(\mathbf{R}^{(+)}) \\
        E^{(-)}, \boldsymbol{\mu}^{(-)} &= \text{SCF}(\mathbf{R}^{(-)})
    \end{align}

    \begin{minted}{python}
    mf_pos = scf.RHF(mol_pos); mf_pos.kernel()
    dip_pos = mf_pos.dip_moment(unit='Bohr')

    mf_neg = scf.RHF(mol_neg); mf_neg.kernel()
    dip_neg = mf_neg.dip_moment(unit='Bohr')
    \end{minted}

    \item \textbf{Обчислення похідної дипольного моменту:}

    Чисельна похідна за нормальною координатою $Q_k$:
    \begin{equation}
        \frac{\partial \boldsymbol{\mu}}{\partial Q_k} = \frac{\boldsymbol{\mu}^{(+)} - \boldsymbol{\mu}^{(-)}}{2h}
    \end{equation}

    \begin{minted}{python}
    dip_derivative = (dip_pos - dip_neg) / (2 * h)
    \end{minted}

    \item \textbf{Обчислення інтегральної інтенсивності:}

    Інтенсивність ІЧ-переходу пропорційна квадрату норми похідної дипольного моменту:
    \begin{equation}
        I_k = 974.8638 \times \left|\frac{\partial \boldsymbol{\mu}}{\partial Q_k}\right|^2_{\text{a.u.}}
        \quad \text{[км/моль]}
    \end{equation}

    \begin{minted}{python}
    intensity = np.linalg.norm(dip_derivative)**2 * 974.8638
    intensities.append(intensity)
    \end{minted}
\end{enumerate}



Хоча метод простий, він вимагає виконання двох \texttt{SCF}-розрахунків на кожну моду
і може бути чутливим до вибору~$\delta$.

%---------------------------------------------------------
\subsection{Аналітичний підхід у \texttt{PySCF}}
%---------------------------------------------------------
Модуль \inlinecode{pyscf.prop.freq} реалізує аналітичний обрахунок похідних
дипольного моменту по атомних координатах $\partial\boldsymbol{\mu}/\partial \mathbf{R}$.

\paragraph{Перехід до нормальних координат.}
Похідна дипольного моменту за нормальною координатою $Q_i$ виражається через
матрицю нормальних мод $\mathbf{L}$:
\[
\frac{\partial \vect{\mu}}{\partial Q_i}
= \sum_{A=1}^{N_{\text{атм}}}
    \frac{\partial \vect{\mu}}{\partial \vect{R}_{A}} \cdot \frac{\partial \vect{R}_A}{\partial Q_i}
= \sum_{A=1}^{N_{\text{атм}}}
    \frac{\partial \vect{\mu}}{\partial \vect{R}_{A}} \cdot \vect{L}_A^i,
\]
де $\vect L_A^i = (L_{Ax}^i, L_{Ay}^i, L_{Az}^i)$ --- вектор нормальної моди для атома $A$.

\paragraph{Структура тензора похідних.}
У \texttt{PySCF} тензор похідних \inlinecode{vib.de} має розмірність $(N_{\text{атм}}, N_{\text{атм}}, 3, 3)$
і зберігає похідні компонент дипольного моменту за координатами атомів:
\[
(\texttt{vib.de})_{ABaj} = \frac{\partial \mu_j^{(A)}}{\partial R_{Ba}},
\]

Діагональні елементи ($A=B$) відповідають похідній внеску атома в дипольний момент по його власних координатах:
\[
D_{Aaj} = (\texttt{vib.de})_{AAaj}
= \frac{\partial \mu_j}{\partial R_{Aa}}.
\]

\begin{minted}{python}
# Діагональний тензор: (N_atm, 3, 3)
de_diag = np.einsum('AAaj->Aaj', vib.de)
\end{minted}

Для кожного атома $A$ отримуємо матрицю Якобі дипольного моменту:
\[
\mathbf{D}_A =
\begin{pmatrix}
\partial \mu_x/\partial x_A & \partial \mu_y/\partial x_A & \partial \mu_z/\partial x_A \\
\partial \mu_x/\partial y_A & \partial \mu_y/\partial y_A & \partial \mu_z/\partial y_A \\
\partial \mu_x/\partial z_A & \partial \mu_y/\partial z_A & \partial \mu_z/\partial z_A
\end{pmatrix}.
\]

\paragraph{Похідні за нормальними координатами.}
Матриця похідних дипольного моменту за всіма нормальними модами:
\[
\frac{\partial \mu_j}{\partial Q_i}
= \sum_{A=1}^{N_{\text{атм}}} \sum_{a=x,y,z} L_{Aa}^i \, D_{Aaj}.
\]

\begin{minted}{python}
# Тензорна згортка: (N_modes, 3)
dipole_derivs = np.einsum('iAa,Aaj->ij', modes, de_diag)
\end{minted}

\paragraph{Інтенсивності ІЧ-переходів.}
Інтенсивність $i$-го коливального переходу:
\[
I_i = 974.8638 \sum_{j=x,y,z} \left( \frac{\partial \mu_j}{\partial Q_i} \right)^2
= 974.8638 \, \left\| \frac{\partial \boldsymbol{\mu}}{\partial Q_i} \right\|^2
\]

\begin{minted}{python}
# Інтенсивності для всіх мод
ir_intensities = 974.8638 * np.sum(dipole_derivs**2, axis=1)
\end{minted}

І весь код разом:

\begin{minted}{python}
import numpy as np
from pyscf import gto, scf
from pyscf.prop import freq

mol = gto.M(atom='O 0 0 0; H 0 -0.757 0.587; H 0 0.757 0.587', basis='6-31g')
mf = scf.RHF(mol).run()

vib = freq.rhf.Frequency(mf)
freqs, modes = vib.kernel()
natm = mol.natm

# Діагональні елементи
de_diag = np.einsum('AAaj->Aaj', vib.de)  # (natm, 3, 3)

# Всі дипольні похідні одразу
dipole_derivs = np.einsum('iAa,Aaj->ij', modes, de_diag)

# Всі інтенсивності одразу
ir_intensities = 974.8638 * np.sum(dipole_derivs**2, axis=1)

# Вивід
print("\n" + "="*70)
print("  ІЧ СПЕКТР H₂O (PySCF, оптимізовано)")
print("="*70)
print(f"{'Мода':<8}{'Частота (см⁻¹)':>18}{'Інтенсивність (km/mol)':>28}")
print("-"*70)
for i in range(6, len(freqs)):
    print(f"{i-5:<8d}{freqs[i]:>18.2f}{ir_intensities[i]:>28.2f}")
print("="*70)
\end{minted}

На відміну від чисельного підходу, цей метод:
\begin{itemize}
    \item не потребує перерахунку енергії при зсуві атомів;
    \item базується на аналітичних градієнтах і похідних;
    \item забезпечує вищу точність і швидкість.
\end{itemize}


\subsubsection{Результати для \ce{H2O}}

\paragraph{Нормальні моди води.}

Експериментальні дані

\begin{table}[h]
\centering
\small
\caption{Експериментальні дані коливань молекули H$_2$O (джерела: \url{https://cccbdb.nist.gov/expvibs2x.asp}).}
\label{tab:exp_vib_H2O}
\begin{tblr}{
colspec = {Q[c,m] Q[l,m,] Q[c,m] Q[c,m] Q[c,m] X[l,m]},
row{1} = {font=\bfseries, c},
hline{1,2,Z} = {-}{0.08em},
}
Мода & Симетрія & {Частота\\ гармонічна (см$^{-1}$)} & {Інтенсивність\\ (км моль$^{-1}$)} \\
1 & \ce{H-O-H} вигин &  1648.5 & 62.5  \\
2 & \ce{O-H} сим. розтяг &  3832.2 & 2.93 \\
3 & \ce{O-H} антисим. розтяг  & 3942.5 & 41.7  \\
\end{tblr}
\end{table}


Молекула H$_2$O має $3N - 6 = 3$ коливальні моди:

\begin{center}
\begin{tblr}{
colspec = {Q[c,m] Q[c,m] Q[c,m] Q[c,m]},
row{1} = {font=\bfseries, c, m},
hline{1, Z} = {solid},
}
Мода & $\nu$ (см$^{-1}$) & $I$ (км/моль) & Тип \\
1 & 1828.9 & 1.4 & Слабка (деформаційна) \\
2 & 3906.0 & 102.3 & Сильна (симетричне валентне) \\
3 & 4001.1 & 138.6 & Сильна (асиметричне валентне) \\
\end{tblr}
\end{center}

\smallskip
\textit{Примітка.} Усі три моди молекули \ce{H2O} є ІЧ-активними
(група симетрії $C_{2v}$). Найсильніше поглинання відповідає
валентним коливанням, що зумовлюють інтенсивну ІЧ-смугу пари води
в атмосфері (важливий внесок у парниковий ефект).

\subsection*{Резюме}

У цьому розділі розглянуто теоретичні основи та практичні методи розрахунку
інфрачервоних спектрів молекул.

\paragraph{Ключові положення:}
\begin{itemize}
    \item ІЧ-інтенсивність коливальної моди визначається квадратом похідної
    дипольного моменту за нормальною координатою:
    $I_i = 974.8638 \left|\frac{\partial\boldsymbol{\mu}}{\partial Q_i}\right|^2$ [км/моль].

    \item Інтенсивність відображає \emph{силу взаємодії} молекули з ІЧ-випромінюванням:
    чим більша зміна дипольного моменту при коливанні, тим інтенсивніша смуга поглинання.
    Класифікація: сильні смуги ($I > 100$~км/моль), середні ($10 < I \leq 100$),
    слабкі ($I \leq 10$~км/моль).

    \item Два методи обчислення похідних:
    \begin{itemize}
        \item \textit{Чисельний} --- простий, але вимагає додаткових SCF-розрахунків
        для зміщених геометрій та чутливий до вибору кроку $h$;
        \item \textit{Аналітичний} --- ефективніший і точніший, базується на аналітичних
        градієнтах густини заряду. Реалізований у \texttt{PySCF} як для HF,
        так і для DFT-методів.
    \end{itemize}

    \item Для молекули H$_2$O розраховані частоти ($\sim 1829$, $3906$, $4001$~см$^{-1}$)
    якісно узгоджуються з експериментом ($1648$, $3832$, $3943$~см$^{-1}$),
    а систематичне завищення пояснюється гармонічним наближенням і вибором базису.

    \item Інтенсивності ($I_1 = 1.4$, $I_2 = 102.3$, $I_3 = 138.6$~км/моль)
    демонструють, що найінтенсивніші смуги відповідають валентним коливанням \ce{O-H} зв'язків,
    тоді як деформаційна мода має слабку інтенсивність. Експериментальні значення
    ($62.5$, $2.93$, $41.7$~км/моль) показують іншу картину через ангармонічні ефекти
    та обмеження гармонічного наближення.

    \item Висока інтенсивність валентних коливань робить пару води потужним поглиначем
    ІЧ-випромінювання в атмосфері, що має значення для клімату та парникового ефекту.
\end{itemize}

\paragraph{Практичне значення:}
Розрахунки ІЧ-спектрів є важливим інструментом для:
\begin{itemize}
    \item ідентифікації функціональних груп та структури молекул;
    \item інтерпретації експериментальних спектрів та передбачення відносних інтенсивностей;
    \item передбачення спектральних характеристик нових сполук;
    \item дослідження міжмолекулярних взаємодій та динаміки;
    \item кількісного аналізу концентрацій речовин за допомогою закону Бугера-Ламберта-Бера.
\end{itemize}



%%% --------------------------------------------------------
\section{Раманівський спектр}
%% --------------------------------------------------------

\textbf{Раман-спектроскопія} досліджує не поглинання, а \emph{розсіювання} світла на молекулах.
Коли фотон із частотою $\nu_0$ падає на молекулу, більшість випромінюється назад із тією ж частотою — це \emph{Релеївське розсіювання}.
Однак невелика частина розсіюється з іншою частотою, зсуненою на коливальну величину $\nu_k$:
\[
\nu_{\text{розс}} = \nu_0 \pm \nu_k,
\]
де «мінус» відповідає \emph{лінії Стокса}, а «плюс» — \emph{анти-Стоксу}.

\textbf{Фізичний зміст.}
Світлове поле $\mathbf{E}$ індукує у молекулі миттєвий дипольний момент:
\[
\boldsymbol{\mu}_{\text{інд}} = \boldsymbol{\alpha}\,\mathbf{E},
\]
де $\boldsymbol{\alpha}$ — тензор поляризовності, що описує, як електронна хмара реагує на поле.
Під час коливання молекули параметр $\boldsymbol{\alpha}$ змінюється.
Саме ця зміна створює компоненту поля на частоті $\nu_0 \pm \nu_k$ — тобто Раманівське розсіювання.

\textbf{Умова активності.}
Коливальна мода $Q_k$ є Раман-активною, якщо при русі ядер поляризовність змінюється:
\[
\frac{\partial \boldsymbol{\alpha}}{\partial Q_k} \neq 0.
\]
Інтенсивність розсіювання для моди $k$ пропорційна квадрату цієї похідної:
\[
I_k^{\text{Раман}} \propto
\left|
\frac{\partial \alpha_{ij}}{\partial Q_k}
\right|^2,
\]
де індекси $i,j=x,y,z$ позначають компоненти тензора поляризовності.

\vspace{1em}
\textbf{Як моделюють Раман-спектр на практиці:}
\begin{enumerate}
    \item \textbf{Коливальні моди.}
    З оптимізованої геометрії обчислюють гесіан $\mathbf{H}$ та отримують нормальні частоти $\omega_k$ і вектори мод $\mathbf{L}_k$.

    \item \textbf{Симетрія.}
    За допомогою теорії груп визначають, які моди Раман-активні (для цього аналізують перетворення тензора $\boldsymbol{\alpha}$ у даній точковій групі).

    Для молекул із центром інверсії (наприклад, \ce{CO2}) діє \emph{правило взаємовиключності}:
    \begin{itemize}
        \item парні моди ($g$) — \textbf{Раман-активні}, але ІЧ-неактивні;
        \item непарні моди ($u$) — \textbf{ІЧ-активні}, але Раман-неактивні.
    \end{itemize}

    \item \textbf{Похідна поляризовності.}
    Для кожної моди обчислюють зміну тензора $\boldsymbol{\alpha}$ при малих зміщеннях $\pm\Delta Q_k$ вздовж нормального вектора:
    \[
    \left. \frac{\partial \alpha_{ij}}{\partial Q_k} \right|_0
    \approx
    \frac{\alpha_{ij}(+\Delta Q_k) - \alpha_{ij}(-\Delta Q_k)}{2\Delta Q_k}.
    \]

    \item \textbf{Інтенсивність лінії.}
    Вводять два інваріанти тензора похідних:
    \[
    \bar{\alpha}_k = \tfrac{1}{3} \mathrm{tr}\!\left(\frac{\partial \boldsymbol{\alpha}}{\partial Q_k}\right),
    \qquad
    \gamma_k^2 = \tfrac{1}{2}\!\!\sum_{i<j}\!\!\left[\!
    \left(\frac{\partial \alpha_{ii}}{\partial Q_k} - \frac{\partial \alpha_{jj}}{\partial Q_k}\right)^{\!2}
    + 6\!\left(\frac{\partial \alpha_{ij}}{\partial Q_k}\right)^{\!2}\right]\!.
    \]
    Тоді відносна Раман-активність:
    \[
    I_k = 45\,\bar{\alpha}_k^2 + 7\,\gamma_k^2.
    \]

    \item \textbf{Побудова спектра.}
    Кожній моді з частотою $\nu_k$ приписують інтенсивність $I_k$ та накладають профіль лінії (зазвичай Лоренців, ширина $\Gamma \approx 5~\text{см}^{-1}$):
    \[
    I(\nu) = \sum_k I_k \, g(\nu - \nu_k).
    \]
\end{enumerate}

\textbf{Резюме.}
Якщо при коливанні молекули змінюється її поляризовність — мода видно в Раман-спектрі.
Якщо ж змінюється дипольний момент — мода проявляється в ІЧ-спектрі.
Тому ІЧ та Раман часто дають взаємодоповнюючу інформацію про коливальні стани молекули.

%% --------------------------------------------------------
\subsubsection{Раманівського спектр молекули \ce{H2O}}
%% --------------------------------------------------------


В програмі \inputcode{h2o_raman_activity.py} наведено розрахунок раманівського спектру молекули \ce{H2O}.

\textbf{Типові активності для \ce{H2O}:}
\begin{center}
    \begin{tabular}{lcc}
        \toprule
        Мода    & $\nu$ (см$^{-1}$) & Раман-активність (Å$^4$/amu) \\
        \midrule
        $\nu_1$ & 3657              & 1.8                          \\
        $\nu_2$ & 1595              & 0.4                          \\
        $\nu_3$ & 3756              & 3.2                          \\
        \bottomrule
    \end{tabular}
\end{center}


\textbf{Принцип взаємного виключення:}

Для молекул з центром інверсії (наприклад, \ce{CO2}):
\begin{itemize}
    \item Парні моди ($g$): Раман активні, ІЧ неактивні
    \item Непарні моди ($u$): ІЧ активні, Раман неактивні
\end{itemize}

%\subsubsection{Раман спектр бензену}
%
%\inputcode{benzene_raman_spectrum.py}
%
%\textbf{Вибрані моди бензену \ce{C6H6} (D$_{6h}$):}
%
%\begin{center}
%    \begin{tblr}{
%        colspec = {X[c] X[c] X[r] X[c] X[c]},
%        row{1} = {font=\bfseries},
%        hline{1,2,Z} = {solid},
%            }
%        Мода       & Симетрія & Частота     & ІЧ        & Раман            \\
%                   &          & (см$^{-1}$) &                              \\
%        $\nu_1$    & A$_{1g}$ & 993         & Неактивна & Активна          \\
%        $\nu_2$    & A$_{1g}$ & 3062        & Неактивна & Активна (сильна) \\
%        $\nu_6$    & E$_{2g}$ & 608         & Неактивна & Активна          \\
%        $\nu_{18}$ & E$_{1u}$ & 1038        & Активна   & Неактивна        \\
%        $\nu_{19}$ & E$_{1u}$ & 3080        & Активна   & Неактивна        \\
%    \end{tblr}
%\end{center}
%
%\textit{Розрахунок: B3LYP/6-311+G(2d,p)}
%
%%---------------------------------------------------------
%\begin{figure}[h!]\centering
%%    \includegraphics[width=0.85\linewidth]{\currfiledir/benzene_raman_spectrum.pdf}
%    \caption{Раман спектр бензену демонструє принцип взаємного виключення.
%        Симетричні моди ($g$) активні в Раман, але неактивні в ІЧ.}
%    \label{pic:benzene_raman}
%\end{figure}
%%---------------------------------------------------------
%
%\textbf{Характерні Раман смуги:}
%\begin{itemize}
%    \item 993 см$^{-1}$: дихальна мода кільця (ring breathing)
%    \item 3062 см$^{-1}$: симетричне валентне \ce{C-H}
%    \item 608 см$^{-1}$: деформація кільця
%\end{itemize}
%
%\subsubsection{Порівняння ІЧ та Раман}
%
%\begin{center}
%    \begin{tblr}{
%        colspec = {X[l,2] X[l,3] X[l,3]},
%        row{1} = {font=\bfseries},
%        hline{1,2,Z} = {solid},
%            }
%        Властивість               & ІЧ спектроскопія         & Раман спектроскопія         \\
%        Правило відбору           & $\frac{d\mu}{dQ} \neq 0$ & $\frac{d\alpha}{dQ} \neq 0$ \\
%        Джерело                   & ІЧ лампа                 & Лазер (видиме/УФ)           \\
%        Зразок                    & Розчин, плівка, KBr      & Розчин, кристал             \\
%        Вода                      & Сильно поглинає          & Слабкий сигнал              \\
%        Скло                      & Непрозоре                & Прозоре                     \\
%        Полярні групи             & Сильний сигнал           & Слабкий                     \\
%        Неполярні                 & Слабкий                  & Сильний сигнал              \\
%        \ce{C=C}, \ce{C\bond{3}C} & Слабкі/неактивні         & Сильні                      \\
%        \ce{C=O}, \ce{O-H}        & Дуже сильні              & Слабкі                      \\
%    \end{tblr}
%\end{center}
%% ========================================================
\section{Термохімія та нульова енергія}
%% ========================================================


Розглянемо, як із результатів квантово-хімічного
розрахунку (зокрема з гесіана) отримують термохімічні параметри:
нульову коливальну енергію, термічні поправки, ентропію та вільну енергію.
Ці параметри \emph{відносяться до газової фази} та
розраховуються в рамках \emph{моделі ідеального газу}.

У квантово-хімічних програмах, таких як \texttt{PySCF},
припускається, що молекула є ізольованою частинкою у газовій фазі,
а її внутрішня енергія складається з кількох незалежних внесків:
\[
E_{\text{total}} = E_{\text{elec}} + E_{\text{vib}} + E_{\text{rot}} + E_{\text{trans}}.
\]
Тут $E_{\text{elec}}$ --- \emph{електронна енергія системи},
отримана з розв’язку електронного рівняння Шредінґера
при фіксованих положеннях ядер (наближення Борна–Оппенгеймера).
Залежно від рівня теорії, цей внесок може бути визначений
у межах методу Хартрі–Фока (HF), теорії функціонала густини (DFT)
або більш точних пост-ХФ підходів

Для оцінки цих ядерних внесків використовується
\emph{наближення жорсткого ротатора} (для обертання молекули)
та \emph{гармонічного осцилятора} (для коливань атомів).
Відповідно, термохімічні функції молекул —
енергія, ентальпія, ентропія та вільна енергія Гіббса —
залежать не лише від електронної будови, але й від термічних ефектів,
пов’язаних із поступальним, обертальним і коливальним рухом ядер.



%% --------------------------------------------------------
\subsection{Нульова коливальна енергія (ZPE)}
%% --------------------------------------------------------

Навіть при абсолютному нулі температури ($T=0$~K) молекула не є нерухомою.
Згідно з принципом невизначеності Гайзенберґа, ядра не можуть одночасно
мати точно визначені координати й імпульси, тому вони завжди здійснюють
квантові коливання навколо рівноважного положення.

Для гармонічного осцилятора енергії рівнів задаються:
\[
E_v = \left(n + \frac{1}{2}\right)\hbar\omega, \qquad n = 0, 1, 2, \ldots
\]
Тобто навіть у найнижчому стані ($n=0$) коливальна енергія дорівнює
$\tfrac{1}{2}\hbar\omega$.
Для багатомодової молекули (з $3N-6$ нормальних мод)
нульова коливальна енергія (Zero Point Energy, ZPE) становить:
\[
E_{\text{ZPE}} = \sum_{i=1}^{3N-6} \frac{1}{2}\hbar\omega_i.
\]

Знання ZPE необхідне для точних термохімічних розрахунків (зокрема енергій реакцій), також вона изначає \emph{ізотопні ефекти}, впливає на швидкість \emph{тунелювання} через потенційні бар’єри. Для невеликих органічних молекул становить зазвичай $5$–$15$~ккал/моль .


%% --------------------------------------------------------
\subsection{Термохімічні поправки}
%% --------------------------------------------------------

При підвищенні температури ядра можуть займати збуджені коливальні,
обертальні й поступальні стани. Це змінює середню енергію системи,
її ентальпію, ентропію та вільну енергію Гіббса.

Для кожного ступеня вільності маємо:

\begin{equation*}
    E_{\text{vib}}(T) = \sum_i \frac{\hbar\omega_i}{\exp(\hbar\omega_i/k_BT) - 1},
\end{equation*}
\begin{equation*}
    E_{\text{rot}}(T) =
    \begin{cases}
        k_BT, & \text{для лінійної молекули;} \\
        \frac{3}{2}k_BT, & \text{для нелінійної.}
    \end{cases}
\end{equation*}
\begin{equation*}
    E_{\text{trans}}(T) = \frac{3}{2}k_BT.
\end{equation*}

Тоді повна енергія системи:
\begin{equation*}
    E_{\text{total}}(T) = E_{\text{elec}} + E_{\text{ZPE}} + E_{\text{vib}}(T)
    + E_{\text{rot}}(T) + E_{\text{trans}}(T).
\end{equation*}

Ентальпія визначається як:
\begin{equation*}
    H(T) = E_{\text{total}}(T) + RT.
\end{equation*}

%% --------------------------------------------------------
\subsection{Ентропія молекули}
%% --------------------------------------------------------

Ентропія --- це міра невпорядкованості системи або кількості доступних мікростанів.
Для ідеального газу вона має такі складові:
\[
    S(T) = S_{\text{trans}} + S_{\text{rot}} + S_{\text{vib}} + S_{\text{elec}}.
\]

\paragraph{1. Поступальна ентропія:}
\[
    S_{\text{trans}} = R\left[
        \ln\!\left(\frac{(2\pi m k_BT)^{3/2}k_BT}{h^3 P}\right) + \frac{5}{2}
    \right],
\]
де $m$ --- маса молекули, $P$ --- тиск.

\paragraph{2. Обертальна ентропія:}
\[
    S_{\text{rot}} = R\left[
        \ln\!\left(
        \frac{\sqrt{\pi}T^{3/2}}{\sigma}
        \left(\frac{8\pi^2k_B}{h^2}\right)^{3/2}
        \sqrt{I_A I_B I_C}
        \right) + \frac{3}{2}
    \right],
\]
де $\sigma$ --- число симетрійних перестановок, $I_A,I_B,I_C$ --- головні моменти інерції.

\paragraph{3. Коливальна ентропія:}
\[
    S_{\text{vib}} = R \sum_i
        \left[
            \frac{\theta_i/T}{\exp(\theta_i/T) - 1}
            - \ln\!\left(1 - e^{-\theta_i/T}\right)
        \right],
\]
де $\theta_i = h\nu_i/k_B$ --- характеристична температура моди.

\paragraph{4. Електронна ентропія:}
Зазвичай нехтують, бо основний стан є єдиним при кімнатних температурах.


%% --------------------------------------------------------
\subsubsection{Вільна енергія Гіббса}
%% --------------------------------------------------------

\[
    G(T) = H(T) - TS(T)
\]

Вона визначає термодинамічну стабільність сполук
і напрямок хімічних реакцій.

%% --------------------------------------------------------
\subsection{Як це реалізовано в \texttt{PySCF}}
%% --------------------------------------------------------
У пакеті \texttt{PySCF} повний термохімічний аналіз молекули базується
на енергії електронів $E_{\text{elec}}$, до якої послідовно додаються
внески нульових коливань, теплові поправки та ентропійні терміни. Термохімічні величини обчислюються
на основі гармонічних частот, отриманих із гессіана
(матриці других похідних енергії за координатами атомів),
і додаються до електронної енергії $E_{\text{elec}}$,
отриманої з розв’язку рівнянь Хартрі–Фока або DFT.

Далі програма:
\begin{enumerate}
    \item Визначає нормальні частоти $\nu_i$ після масового масштабування гессіана.
    \item Обчислює нульову енергію коливань $E_{\text{ZPE}} = \tfrac{1}{2}\sum h\nu_i$.
    \item Розраховує термічні внески $E_{\text{vib}}(T), E_{\text{rot}}(T), E_{\text{trans}}(T)$
          за статистично-механічними формулами (як наведено вище).
    \item Обчислює ентальпію $H(T)$, ентропію $S(T)$ та вільну енергію $G(T)$.
\end{enumerate}

Всі ці розрахунки виконуються у межах
\emph{наближення гармонічного осцилятора та жорсткого ротатора},
що забезпечує добру точність для малих і середніх молекул.


%% --------------------------------------------------------
\subsection{Термохімічний приклад: \ce{H2O}}
%% --------------------------------------------------------

\inputcode[mod=inline]{h2o_thermochemistry.py}


У цій програмі основна робота термохімічного аналізу була зведена до двох ключових функцій з модуля \href{https://pyscf.org/_modules/pyscf/hessian/thermo.html}{\texttt{pysch.hessian.thermo}}:

\begin{enumerate}
  \item \inlinecode{harmonic\_analysis(mol, hess)} --- виконує гармонічний аналіз молекули на основі гесіану.
        Вона обчислює частоти нормальних коливань, нормалізовані моди, редуковані маси та нульову коливальну енергію (ZPE).
        Ця функція також враховує трансляційні та обертальні ступені свободи та дозволяє виключати їх при необхідності.
  \item \inlinecode{thermo(model, freq, temperature, pressure)} --- використовує результати гармонічного аналізу для обчислення термодинамічних величин:
        ентропії, теплоємності ($C_V$, $C_P$), внутрішньої енергії, ентальпії та вільної енергії Гіббса.
        Вона враховує внески електронної, поступальної, обертальної та коливальної частин.
\end{enumerate}

Таким чином, за рахунок цих двох функцій програма зводить складні квантово-хімічні розрахунки та гармонічну апроксимацію в зручний формат для термохімічного аналізу.
Вся <<магія>> полягає в тому, що \inlinecode{harmonic\_analysis} формує основні фізичні параметри, а \inlinecode{thermo} перетворює їх у термодинамічні величини, готові для подальшого виводу у таблиці.



%\noindent
%\textit{Коментар:}
%Більшість квантово-хімічних програм (Gaussian, ORCA, PySCF тощо)
%в автоматичному режимі обчислюють ці термохімічні функції на основі
%власних частот, отриманих із гесіану при рівноважній геометрії.
%Тому результати часто подаються у вигляді:
%\[
%E_{\text{elec}}, \quad E_{\text{elec}} + \text{ZPE}, \quad H(298), \quad G(298).
%\]
%Саме ці значення використовують у подальших термодинамічних розрахунках,
%наприклад, при оцінюванні енергії реакції чи рівноважних констант.



%%% --------------------------------------------------------
%\subsubsection{Термохімічний розрахунок для \ce{CH4}}
%%% --------------------------------------------------------
%
%
%\inputcode{ch4_thermochemistry.py}
%
%\textbf{Результати для \ce{CH4} при $298.15$~K:}
%
%\begin{center}
%    \begin{tblr}{
%        colspec = {X[l,3] X[r,2] X[l,2]},
%        row{1} = {font=\bfseries},
%        hline{1,2,Z} = {solid},
%            }
%        Величина              & Значення & Одиниці      \\
%        ZPE                   & 28.03    & ккал/моль    \\
%        $E_{\text{vib}}(T)$   & 0.10     & ккал/моль    \\
%        $E_{\text{rot}}(T)$   & 0.89     & ккал/моль    \\
%        $E_{\text{trans}}(T)$ & 0.89     & ккал/моль    \\
%        $H(T) - H(0)$         & 2.48     & ккал/моль    \\
%        $S(T)$                & 44.5     & кал/(моль·K) \\
%        $C_V$                 & 6.0      & кал/(моль·K) \\
%    \end{tblr}
%\end{center}
%
%\textit{Розрахунок: B3LYP/6-311+G(2d,p)}
%
%\textbf{Порівняння з експериментом:}
%\begin{itemize}
%    \item ZPE (експ.): 27.8 ккал/моль --- відмінна згода
%    \item $S_{298}$ (експ.): 44.5 кал/(моль·K) --- ідеальна згода
%    \item Термохімічні поправки надійні для DFT
%\end{itemize}

%%% --------------------------------------------------------
%\subsection{Ізотопні ефекти}
%%% --------------------------------------------------------
%
%Заміна ізотопу змінює частоти через зміну маси:
%\[
%    \frac{\omega_{\text{D}}}{\omega_{\text{H}}} \approx \sqrt{\frac{m_{\text{H}}}{m_{\text{D}}}} = \sqrt{\frac{1}{2}} \approx 0.707
%\]

%\subsubsection{Ізотопний ефект у \ce{H2O}/\ce{D2O}}
%
%%\inputcode{h2o_d2o_isotope_effect.py}
%
%\textbf{Порівняння \ce{H2O} та \ce{D2O}:}
%
%\begin{center}
%    \begin{tblr}{
%        colspec = {X[l,3] X[r,2] X[r,2] X[r,2]},
%        row{1} = {font=\bfseries},
%        hline{1,2,Z} = {solid},
%            }
%        Мода                & \ce{H2O}    & \ce{D2O}    & Співвідношення \\
%                            & (см$^{-1}$) & (см$^{-1}$) &                \\
%        Симетричне валентне & 3657        & 2671        & 0.730          \\
%        Деформаційне        & 1595        & 1178        & 0.738          \\
%        Антисим. валентне   & 3756        & 2788        & 0.742          \\
%    \end{tblr}
%\end{center}
%
%\textit{Експериментальні дані}
%
%\textbf{ZPE ефект:}
%\begin{itemize}
%    \item ZPE(\ce{H2O}): 13.26 ккал/моль
%    \item ZPE(\ce{D2O}): 9.66 ккал/моль
%    \item $\Delta$ZPE: 3.60 ккал/моль
%\end{itemize}
%
%\textbf{Наслідки:}
%\begin{itemize}
%    \item Дейтеровані сполуки стабільніші (нижча ZPE)
%    \item Кінетичний ізотопний ефект: $k_{\text{H}}/k_{\text{D}} \approx 7$ для розриву \ce{C-H}
%    \item Використання в механістичних дослідженнях
%\end{itemize}

%% --------------------------------------------------------
\section{Перехідні стани та уявні частоти}
%% --------------------------------------------------------

\subsection{Класифікація стаціонарних точок}

Гессіан дає змогу класифікувати стаціонарні точки потенціальної поверхні енергії (ППЕ) за знаками власних значень:
\[
\lambda_i = \frac{\partial^2 E}{\partial q_i^2}.
\]

\begin{center}
    \begin{tblr}{
        colspec = {X[l,2] X[c,2] X[l,3]},
        row{1} = {font=\bfseries},
        hline{1,2,Z} = {solid},
            }
        Тип                & Уявних частот & Характеристика             \\
        Мінімум            & 0             & Локальний мінімум          \\
        Перехідний стан    & 1             & Сідлова точка 1-го порядку \\
        Сідло 2-го порядку & 2             & Нестабільна структура      \\
    \end{tblr}
\end{center}

\textbf{Уявна частота:} якщо \(\omega^2 < 0\), то \(\omega = i|\omega|\), і відповідний напрямок руху описує нестійку моду.

%% --------------------------------------------------------
\subsection{Інверсія аміаку (\ce{NH3}) як приклад}
%% --------------------------------------------------------

Молекула аміаку має тригранно-пірамідальну структуру симетрії $C_{3v}$: атом Нітрогену розташований над площиною трьох атомів Гідрогену. Дзер\-каль\-но-еквівалентна конфігурація виникає, коли атом N переходить на протилежний бік площини --- процес, відомий як \emph{парасолькова інверсія}.

Якщо ввести координату $h$, що задає відстань атома N від площини H$_3$, то потенціальна енергія має форму подвійного мінімуму:
\[
E(h) \approx E_0 + \tfrac{1}{2} k h^2 + \tfrac{1}{4}\lambda h^4, \qquad \lambda>0.
\]
Пірамідальні структури відповідають точкам $h=\pm h_0$ (локальні мінімуми), а планарна структура ($h=0$) є \textbf{перехідним станом}.



У рівноважній геометрії ($h = h_0$) усі власні значення гесіану додатні, що відповідає стабільному мінімуму:
\[
\lambda_i > 0 \quad \Rightarrow \quad \omega_i = \sqrt{\lambda_i/\mu_i} > 0.
\]

Для планарної конфігурації ($h = 0$) одне власне значення стає від’ємним, і відповідна частота набуває уявного значення:
\[
\omega = i|\omega|, \qquad \omega^2 < 0.
\]
Це свідчить про сідлову точку першого порядку — перехідний стан інверсії аміаку.


У класичному наближенні атом \ce{N} мав би подолати потенціальний бар’єр, однак квантова механіка допускає \emph{тунелювання}.  Енергетичне розщеплення між ними ($\Delta E \approx 0.79~\text{см}^{-1}$) відповідає частоті мікрохвильового переходу, що спостерігається експериментально в \emph{амонійному мазері}.

%% --------------------------------------------------------
\paragraph{Обчислення в \texttt{PySCF}}
%% --------------------------------------------------------


У файлі \inputcode{nh3_inversion.py} наведено повний приклад дослідження
інверсії молекули аміаку. Розрахунок демонструє, як із допомогою
\texttt{PySCF} можна чисельно побудувати профіль потенціальної енергії $E(h)$
та проаналізувати коливальні частоти у мінімумі й у перехідному стані.

Основні етапи скрипта:

\begin{enumerate}
    \item \textbf{Функція \inlinecode{energy\_func(h)}}
    Обчислює енергію системи при фіксованій висоті атома N над площиною трьох атомів H.
    Для кожного значення $h$ створюється новий об'єкт \inlinecode{gto.M},
    виконується \texttt{RHF}-розрахунок, і повертається значення енергії.

    \item \textbf{Функція \inlinecode{analyze\_point(h)}}
    Для заданої конфігурації обчислює:
    \begin{itemize}
        \item енергію $E(h)$;
        \item аналітичний градієнт \inlinecode{mf.Gradients().kernel()};
        \item гессіан \inlinecode{mf.Hessian().kernel()} та відповідні гармонічні частоти;
        \item кількість уявних частот (як ознаку перехідного стану).
    \end{itemize}
    Аналіз частот здійснюється через функцію \inlinecode{thermo.harmonic\_analysis()}.

    \item \textbf{Основний цикл аналізу}
    Задається набір висот $h$ (груба та уточнена сітка) і для кожної точки виконується \inlinecode{analyze\_point(h)}.
    Енергії переводяться у ккал/моль, градієнти контролюють локальність мінімуму.
    Для точок із малим градієнтом виводяться частоти коливань.

    \item \textbf{Пошук перехідного стану}
    Після проходу всіх точок програма вибирає:
    \begin{itemize}
        \item мінімум енергії $E_{\min}$ (стійкі геометрії);
        \item максимум енергії $E_{\text{TS}}$ поблизу $h = 0$;
    \end{itemize}
    Перевіряється, чи в кандидата на TS лише одна уявна частота,
    і якщо так, визначається бар’єр інверсії:
    \[
        \Delta E_{\text{бар}} = (E_{\text{TS}} - E_{\min}) \times 627.51 \text{ ккал/моль}.
    \]

    \item \textbf{Візуалізація профілю потенціальної поверхні}
    Побудова графіка $E(h)$ виконується за допомогою \texttt{matplotlib},
    результат зберігається у файл \texttt{nh3\_inversion.pdf}.
\end{enumerate}

Таким чином, програма дозволяє послідовно отримати:
\begin{itemize}
    \item рівноважні геометрії двох еквівалентних мінімумів (пірамідальних форм NH$_3$);
    \item перехідний стан із планарною структурою ($h = 0$);
    \item уявну частоту інверсії, що відповідає русі атома N через площину H$_3$;
    \item бар’єр інверсії, який кількісно характеризує тунельне розщеплення рівнів.
\end{itemize}
